% Vietnamese translation for OI Public Library
% Source-PDF: 比赛 CONTEST/2021“MINIEYE杯”中国大学生算法设计超级联赛/2021“MINIEYE杯”中国大学生算法设计超级联赛（8）/2021 HDU Multi-University Training Contest 8.pdf
\documentclass[11pt]{article}
\usepackage{oipl-vn}

\title{Đề bài HDU Multi-University Training Contest 8 năm 2021}
\author{}
\date{}

\begin{document}
\maketitle

\origtitle{2021 HDU Multi-University Training Contest 8}
\sourcepath{contest/hdu-2021-training-08-problems.pdf}

\section{Problem A. X-liked Counting}

\noindent
\textbf{Input file:} standard input\\
\textbf{Output file:} standard output\\
\textbf{Time limit:} 6 seconds\\
\textbf{Memory limit:} 128 megabytes

David luôn thích các loại số mới.

Anh định nghĩa một số là $x$-liked nếu một tiền tố hoặc một hậu tố bất kỳ của
nó chia hết cho $x$. Một tiền tố là một đoạn liên tiếp bắt đầu từ vị trí ngoài
cùng bên trái của số và kết thúc ở một vị trí bất kỳ của số. Một hậu tố là một
đoạn liên tiếp bắt đầu ở một vị trí bất kỳ của số và kết thúc ở vị trí ngoài
cùng bên phải của số. Ví dụ, $12$ là một tiền tố của $123$, còn $23$ là một
hậu tố của nó. Bản thân số đó cũng có thể là tiền tố hoặc hậu tố của chính nó.

Bây giờ David muốn biết có bao nhiêu số nguyên $x$-liked trong đoạn $[l,r]$.
Tuy nhiên, vì David không thích chữ số $7$, anh không cho phép chữ số $7$ xuất
hiện trong các số. Vì vậy bạn không cần tính các số như $27,372,774$, v.v.,
vì chúng đều có ít nhất một chữ số bằng $7$.

Hãy cho anh biết có bao nhiêu số nguyên $x$-liked trong đoạn $[l,r]$ mà không
có chữ số nào bằng $7$.

\subsection*{Input}

Dòng đầu tiên chứa một số nguyên $T$ ($1\le T\le 10$), là số bộ dữ liệu.

Với mỗi bộ dữ liệu, có một dòng chứa ba số nguyên không âm $l,r$ và $x$.
\[
  1\le l\le r\le 10^{18},\qquad 1\le x\le 500.
\]

\subsection*{Output}

Với mỗi bộ dữ liệu, in ra một số trên một dòng: số lượng số nguyên $x$-liked
trong đoạn $[l,r]$ mà không có chữ số nào bằng $7$.

\subsection*{Example}

\begin{verbatim}
standard input
1
11 20 4

standard output
5
\end{verbatim}

\subsection*{Note}

Trong ví dụ, có năm số nguyên hợp lệ: $12$ ($12$ chia hết cho $4$), $14$
($4$ chia hết cho $4$), $16$ ($16$ chia hết cho $4$), $18$ ($8$ chia hết cho
$4$), và $20$ ($20$ và $0$ đều chia hết cho $4$).

\section{Problem B. Buying Snacks}

\noindent
\textbf{Input file:} standard input\\
\textbf{Output file:} standard output\\
\textbf{Time limit:} 3 seconds\\
\textbf{Memory limit:} 256 megabytes

Là một người sành ăn, Diana rất thích ăn đồ ăn vặt. Một ngày nọ, toàn bộ đồ
ăn vặt của cô bị Bella tịch thu, nên Diana quyết định mua thêm ở một cửa hàng
đồ ăn vặt.

Trong cửa hàng có tổng cộng $n$ loại đồ ăn vặt khác nhau, được đánh số
$1,2,\ldots,n$. Mỗi loại có hai cỡ gói: gói lớn và gói nhỏ. Diana muốn nếm
càng nhiều loại đồ ăn vặt khác nhau càng tốt, nên với mỗi loại cô sẽ mua
nhiều nhất một gói. Với hai loại đồ ăn vặt kề nhau bất kỳ, không phụ thuộc cỡ
gói, Diana có thể chọn mua chúng thành một combo để được giảm giá so với mua
riêng lẻ. Để đơn giản hóa bài toán, giả sử mọi gói nhỏ có giá $1$ đồng, mọi
gói lớn có giá $2$ đồng, và mỗi lần giảm giá giảm $1$ đồng. Diana muốn biết có
bao nhiêu cách khác nhau để mua đồ ăn vặt với đúng $k$ đồng.

Hãy in đáp án modulo $998244353$ cho mỗi $k\in [1,m]$.

Hai cách được xem là giống nhau nếu các loại đồ ăn vặt, cỡ gói và các combo
đều giống nhau.

\textbf{Chú ý:} Với hai loại đồ ăn vặt kề nhau, Diana có thể mua riêng chúng
mà không giảm giá, hoặc mua chúng thành một combo; hai cách này được xem là
khác nhau.

Ví dụ, khi $n=2$, $k=2$, có $5$ cách khác nhau:
\[
  (1b),\ (2b),\ (1s\&2b),\ (1b\&2s),\ (1s)(2s).
\]
Trong đó số là loại đồ ăn vặt, $s$ nghĩa là gói nhỏ, $b$ nghĩa là gói lớn, và
\texttt{\&} nghĩa là combo.

\subsection*{Input}

Dòng đầu tiên chứa một số nguyên $T$ ($T\le 5$), là số bộ dữ liệu. Mỗi bộ dữ
liệu chứa hai số nguyên $n$ và $m$
($1\le n\le 10^9$, $1\le m\le 20000$), lần lượt là số loại đồ ăn vặt và số
đồng tối đa. Dữ liệu bảo đảm rằng với mọi bộ dữ liệu,
\[
  \sum m \le 30000.
\]

\subsection*{Output}

Với mỗi bộ dữ liệu, in một dòng gồm $m$ số nguyên, biểu thị các đáp án.

\subsection*{Example}

\begin{verbatim}
standard input
2
2 3
5 4

standard output
3 5 3
9 38 97 166
\end{verbatim}

\section{Problem C. Ink on paper}

\noindent
\textbf{Input file:} standard input\\
\textbf{Output file:} standard output\\
\textbf{Time limit:} 5 seconds\\
\textbf{Memory limit:} 512 megabytes

Bob vô tình làm đổ vài giọt mực lên giấy. Vị trí ban đầu của giọt mực thứ $i$
là $(x_i,y_i)$; giọt mực lan ra ngoài với tốc độ $0.5$ xăng-ti-mét mỗi giây,
tạo thành một hình tròn. Bob tò mò muốn biết sau bao lâu thì tất cả các vết
mực trở nên liên thông. Để thuận tiện cho việc xuất kết quả, hãy in bình
phương của thời gian đó.

\subsection*{Input}

Dòng đầu tiên chứa một số nguyên $T$ ($1\le T\le 5$), là số bộ dữ liệu. Với
mỗi bộ dữ liệu, dòng đầu tiên chứa một số nguyên $n$ ($2\le n\le 5000$), là
số giọt mực trên giấy. Mỗi dòng trong $n$ dòng tiếp theo chứa hai số nguyên
$(x_i,y_i)$ ($|x_i|\le 10^9$, $|y_i|\le 10^9$), là tọa độ $x$ và $y$ của
giọt mực.

\subsection*{Output}

Với mỗi bộ dữ liệu, in một dòng chứa một số thập phân, là đáp án.

\subsection*{Example}

\begin{verbatim}
standard input
2
3
0 0
1 1
0 1
5
1 1
4 5
1 4
2 6
3 10

standard output
1
17
\end{verbatim}

\section{Problem D. Counting Stars}

\noindent
\textbf{Input file:} standard input\\
\textbf{Output file:} standard output\\
\textbf{Time limit:} 1 second\\
\textbf{Memory limit:} 128 megabytes

Là một cô gái thanh lịch và tài giỏi, Bella thích ngắm sao vào ban đêm, đặc
biệt là sao Bắc Cực. Có tổng cộng $n$ ngôi sao trên bầu trời. Ta định nghĩa
$a_i$ là độ sáng ban đầu của ngôi sao thứ $i$.

Ở giây thứ $t$, độ sáng của một số ngôi sao có thể thay đổi do hoạt động của
sao. Để đơn giản hóa bài toán, ta xấp xỉ chia các hoạt động thành hai loại:
\begin{enumerate}
  \item Với mọi $i\in [l,r]$, độ sáng của ngôi sao thứ $i$ giảm đi
  $a_i\&(-a_i)$.
  \item Với mọi $i\in [l,r]$ sao cho $a_i\ne 0$, độ sáng của ngôi sao thứ $i$
  tăng thêm $2^k$, trong đó $k$ thỏa $2^k\le a_i<2^{k+1}$.
\end{enumerate}

Ngoài ra, Bella có một số truy vấn. Với mỗi truy vấn, cô muốn biết tổng độ
sáng của các ngôi sao trong đoạn $[l,r]$, tức là $\sum_{i=l}^{r} a_i$.

Như chúng ta đều biết, Bella rất thông minh. Vì vậy cô không thể tự giải một
bài khó như thế. Hãy giúp cô giải bài toán và trả lời các truy vấn.

Vì đáp án có thể rất lớn, hãy in đáp án modulo $998244353$.

\subsection*{Input}

Dòng đầu tiên chứa một số nguyên $T$ ($T\le 5$), là số bộ dữ liệu.

Với mỗi bộ dữ liệu:

Dòng đầu tiên chứa một số nguyên $n$ ($1\le n\le 100000$), là số ngôi sao.

Dòng thứ hai chứa $n$ số nguyên $a_1,a_2,\ldots,a_n$ ($1\le a_i\le 10^9$),
là độ sáng ban đầu của các ngôi sao.

Dòng thứ ba chứa một số nguyên $q$ ($1\le q\le 100000$), là số thao tác.

Mỗi dòng trong $q$ dòng tiếp theo chứa ba số nguyên
$\mathit{opt}_i,l_i,r_i$, biểu thị loại thao tác và đoạn thao tác.
\[
\begin{array}{ll}
\mathit{opt}_i=1 & \text{nghĩa là truy vấn của Bella,}\\
\mathit{opt}_i=2 & \text{nghĩa là hoạt động sao loại thứ nhất,}\\
\mathit{opt}_i=3 & \text{nghĩa là hoạt động sao loại thứ hai.}
\end{array}
\]
Dữ liệu bảo đảm rằng với mọi bộ dữ liệu,
\[
  \sum n\le 4\times 10^5,\qquad \sum q\le 4\times 10^5.
\]

\subsection*{Output}

Với mỗi truy vấn của Bella, in đáp án modulo $998244353$ trên một dòng riêng.

\subsection*{Example}

\begin{verbatim}
standard input
1
5
5 2 2 9 7
4
2 1 5
1 1 1
3 1 3
1 2 5

standard output
4
14
\end{verbatim}

\section{Problem E. Separated Number}

\noindent
\textbf{Input file:} standard input\\
\textbf{Output file:} standard output\\
\textbf{Time limit:} 1 second\\
\textbf{Memory limit:} 128 megabytes

Cathy yêu các con số, và gần đây cô thích việc tách các con số.

Một cách tách một số được định nghĩa là chia số đó thành các phần liên tiếp.
Ví dụ, ta có thể gọi $(11)(451)(4)$ là một cách tách số $114514$. Giá trị của
một cách tách là tổng của tất cả các phần sau khi tách; giá trị của
$(11)(451)(4)$ bằng $11+451+4=466$. Nếu một phần có các số không đứng đầu thì
vẫn hợp lệ, nên cách tách $(1)(00)$ của số $100$ cũng là một cách tách hợp lệ.
Bây giờ Cathy có một số $x$ không có số không đứng đầu, và cô muốn biết tổng
giá trị của các cách tách số đó thành không quá $k$ phần. Cô không thật sự
thông minh, nên cô nhờ bạn giúp.

Vì đáp án có thể rất lớn, bạn chỉ cần in đáp án modulo $998244353$.

\subsection*{Input}

Dòng đầu tiên chứa một số $T$ ($1\le T\le 5$), là số bộ dữ liệu.

Với mỗi bộ dữ liệu, có hai dòng.

Dòng đầu tiên chứa một số $k$, là số phần tối đa.

Dòng thứ hai chứa một số $x$, là số được hỏi.

Gọi $n$ là số chữ số của $x$, ta có $1\le n\le 10^6$ và $1\le k\le n$.
Dữ liệu bảo đảm rằng với mọi bộ dữ liệu,
\[
  \sum n\le 10^6.
\]

\subsection*{Output}

Với mỗi bộ dữ liệu, in một số trên một dòng: đáp án modulo $998244353$.

\subsection*{Example}

\begin{verbatim}
standard input
1
3
100

standard output
112
\end{verbatim}

\subsection*{Note}

Trong ví dụ, có $4$ cách tách thành không quá $3$ phần:
$(100)$, $(1)(00)$, $(10)(0)$, $(1)(0)(0)$. Giá trị của chúng lần lượt là
$100$, $1+0=1$, $10+0=10$, $1+0+0=1$, nên đáp án là
$100+1+10+1=112$.

\section{Problem F. GCD Game}

\noindent
\textbf{Input file:} standard input\\
\textbf{Output file:} standard output\\
\textbf{Time limit:} 1 second\\
\textbf{Memory limit:} 128 megabytes

Alice và Bob đang chơi một trò chơi.

Họ luân phiên thao tác. Có $n$ số $a_1,a_2,\ldots,a_n$. Mỗi lần, người chơi
thực hiện $3$ bước:
\begin{enumerate}
  \item Tùy ý chọn một số $a_i$.
  \item Tùy ý chọn một số khác $x$ ($1\le x<a_i$).
  \item Thay số $a_i$ bằng $\gcd(a_i,x)$. Ở đây, $\gcd(u,v)$ là ước chung
  lớn nhất của $u$ và $v$.
\end{enumerate}

Khi một người chơi không thể thực hiện bất kỳ nước đi nào, người đó thua.
Alice đi trước, và cô yêu cầu bạn cho biết ai sẽ thắng nếu cả hai chơi tối ưu.

\subsection*{Input}

Dòng đầu tiên chứa một số $T$ ($1\le T\le 100$), là số bộ dữ liệu.

Với mỗi bộ dữ liệu, có hai dòng.

Dòng đầu tiên chứa một số $n$ ($1\le n\le 10^6$).

Dòng thứ hai chứa $n$ số $a_1,a_2,\ldots,a_n$ ($1\le a_i\le 10^7$).

\subsection*{Output}

Với mỗi bộ dữ liệu, in đáp án trên một dòng.

Nếu Alice sẽ thắng, in \texttt{Alice}; ngược lại in \texttt{Bob}.

\subsection*{Example}

\begin{verbatim}
standard input
2
1
1
1
2

standard output
Bob
Alice
\end{verbatim}

\section{Problem G. A Simple Problem}

\noindent
\textbf{Input file:} standard input\\
\textbf{Output file:} standard output\\
\textbf{Time limit:} 2 seconds\\
\textbf{Memory limit:} 512 megabytes

Bạn có một dãy $A$ độ dài $n$ và một số nguyên dương $k$. Ban đầu, mọi phần tử
trong $A$ đều bằng $0$.

Bây giờ có $q$ thao tác. Các thao tác này được chia thành hai loại:
\[
\begin{array}{ll}
\texttt{1 l r x}: & \forall i\in [l,r],\ A_i=A_i+x,\\[0.2em]
\texttt{2 l r}: & \text{tìm }\displaystyle
  \min_{i=l}^{r-k+1}\left(\max_{j=i}^{i+k-1} a_j\right).
\end{array}
\]

\subsection*{Input}

Dòng đầu tiên chứa một số nguyên $T$ ($T\le 5$), là số bộ dữ liệu. Mỗi bộ dữ
liệu chứa $q+2$ dòng.

Dòng đầu tiên chứa ba số nguyên $n,k$ và $q$
($2\le k\le n\le 5\times 10^8$, $1\le q\le 10^5$).

$q$ dòng tiếp theo mô tả các thao tác thuộc hai loại:
\[
\begin{array}{ll}
\texttt{1 l r x}: & \forall i\in [l,r],\ A_i=A_i+x,\\[0.2em]
\texttt{2 l r}: & \text{tìm }\displaystyle
  \min_{i=l}^{r-k+1}\left(\max_{j=i}^{i+k-1} a_j\right).
\end{array}
\]
Dữ liệu bảo đảm tổng $q$ không vượt quá $2\times 10^5$.

\subsection*{Output}

Với mỗi thao tác loại $2$, in đáp án trên một dòng riêng.

\subsection*{Example}

\begin{verbatim}
standard input
2
5 3 3
1 2 5 2
1 3 4 -1
2 1 4
10 4 10
1 1 6 6
1 3 8 -6
2 2 6
1 4 8 -8
1 4 9 4
1 4 5 -7
2 4 8
1 6 7 8
1 1 3 -2
2 3 7

standard output
2
0
-4
4
\end{verbatim}

\section{Problem H. Square Card}

\noindent
\textbf{Input file:} standard input\\
\textbf{Output file:} standard output\\
\textbf{Time limit:} 1 second\\
\textbf{Memory limit:} 64 megabytes

Eric đang chơi một trò chơi trên một mặt phẳng vô hạn.

Trên mặt phẳng có một vùng tròn bán kính $r_1$ gọi là Vùng Ghi Điểm. Mỗi lần,
Eric ném một tấm thẻ hình vuông cạnh $a$ lên mặt phẳng, rồi tấm thẻ bắt đầu
quay quanh tâm của nó. Nếu tại một thời điểm nào đó tấm thẻ nằm hoàn toàn
nghiêm ngặt bên trong Vùng Ghi Điểm, lượt chơi hiện tại sẽ được tính điểm.

Còn có một vùng tròn khác bán kính $r_2$ gọi là Vùng Thưởng. Nếu lượt chơi
được tính điểm trong tình huống hiện tại, tấm thẻ hình vuông sẽ tiếp tục quay;
nếu tại một thời điểm nào đó nó nằm hoàn toàn nghiêm ngặt bên trong Vùng
Thưởng, Eric sẽ nhận thêm thưởng.

Eric không giỏi trò chơi này, nên ta có thể xem đơn giản rằng anh ném tấm thẻ
đến mọi vị trí với xác suất như nhau.

Bây giờ Eric muốn biết tỉ số giữa xác suất vừa được tính điểm vừa nhận thưởng
và xác suất được tính điểm.

Tọa độ tâm của Vùng Ghi Điểm là $(x_1,y_1)$.

Tọa độ tâm của Vùng Thưởng là $(x_2,y_2)$.

\subsection*{Input}

Dòng đầu tiên chứa một số $T$ ($1\le T\le 20$), là số bộ dữ liệu.

Với mỗi bộ dữ liệu, có ba dòng.

Dòng đầu tiên chứa ba số nguyên $r_1,x_1,y_1$
($-1000\le x_1,y_1\le 1000$, $0<r_1\le 1000$), là bán kính và vị trí của
Vùng Ghi Điểm.

Dòng thứ hai chứa ba số nguyên $r_2,x_2,y_2$
($-1000\le x_2,y_2\le 1000$, $0<r_2\le 1000$), là bán kính và vị trí của Vùng
Thưởng.

Dòng thứ ba chứa một số nguyên $a$ ($0<a\le 1000$), là độ dài cạnh của tấm
thẻ hình vuông.

Dữ liệu bảo đảm Vùng Ghi Điểm đủ lớn, nên luôn có khả năng lượt chơi được
tính điểm.

\subsection*{Output}

Với mỗi bộ dữ liệu, in một số thập phân $p$ ($0\le p\le 1$) trên một dòng, là
tỉ số giữa xác suất vừa được tính điểm vừa nhận thưởng và xác suất được tính
điểm.

Đáp án của bạn cần được làm tròn tới $6$ chữ số sau dấu thập phân.

\subsection*{Example}

\begin{verbatim}
standard input
1
5 1 2
3 2 1
1

standard output
0.301720
\end{verbatim}

\section{Problem I. Singing Superstar}

\noindent
\textbf{Input file:} standard input\\
\textbf{Output file:} standard output\\
\textbf{Time limit:} 2 seconds\\
\textbf{Memory limit:} 512 megabytes

Ke Ke là một cô gái vui vẻ và nhiệt tình. Cô mơ ước được hát vui vẻ mỗi ngày.
Một ngày nọ, khi đang luyện hát, Ke Ke vô tình phát hiện có một số từ có thể
làm bài hát hay hơn. Cô gọi những từ này là ``SuperstarWords''.

Sau rất nhiều lần thử, cô nhận thấy rằng trong mỗi bài hát có tổng cộng $n$
``SuperstarWords''. Cô muốn biết số lần xuất hiện rời nhau tối đa của từng
``SuperstarWords'' trong bài hát của mình.

\textbf{Chú ý:} một lần xuất hiện nghĩa là SuperstarWord là một xâu con liên
tiếp của bài hát tại một vị trí nào đó.

Ví dụ: số lần xuất hiện rời nhau tối đa của ``aba'' trong ``abababa'' là $2$.

\subsection*{Input}

Dòng đầu tiên chứa một số nguyên $T$ ($T\le 5$), là số bộ dữ liệu.

Dòng đầu tiên của mỗi bộ dữ liệu chứa một xâu $s$ ($1\le |s|\le 10^5$), biểu
thị một bài hát.

Dòng thứ hai của mỗi bộ dữ liệu chứa một số nguyên $n$ ($1\le n\le 10^5$),
biểu thị số lượng SuperstarWords.

$n$ dòng tiếp theo, mỗi dòng chứa một xâu $a_i$ ($1\le |a_i|\le 30$), biểu
thị một SuperstarWord.

Dữ liệu bảo đảm rằng với mọi bộ dữ liệu,
\[
  \sum |s|\le 4\times 10^5,\qquad \sum |a_i|\le 4\times 10^5.
\]

\textbf{Chú ý:} tất cả các xâu trong dữ liệu vào chỉ gồm chữ cái tiếng Anh
thường.

\subsection*{Output}

Với mỗi bộ dữ liệu, in $n$ dòng; mỗi dòng chứa một số nguyên, biểu thị đáp án.

\subsection*{Example}

\begin{verbatim}
standard input
2
abababbaanbanananabanana
3
aba
ababa
nana
nihaoxiexiexiaolongbaozaijian
3
qwer
taihaotinleba
xiexi

standard output
2
1
2
0
0
1
\end{verbatim}

\section{Problem J. Yinyang}

\noindent
\textbf{Input file:} standard input\\
\textbf{Output file:} standard output\\
\textbf{Time limit:} 1 second\\
\textbf{Memory limit:} 512 megabytes

Bạn có một lưới gồm $n$ hàng và $m$ cột; mỗi ô cần được tô đen hoặc trắng.

Ô ở hàng thứ $i$ và cột thứ $j$ được ký hiệu là $(i,j)$.

Hai ô được nối trực tiếp khi và chỉ khi chúng có chung một cạnh.

Hai ô được liên thông khi và chỉ khi chúng được nối trực tiếp, hoặc tồn tại
một ô liên thông với cả hai ô đó.

Một phương án tô màu là tốt khi và chỉ khi nó thỏa mãn ba điều kiện:
\begin{enumerate}
  \item Tất cả ô trắng liên thông.
  \item Tất cả ô đen liên thông.
  \item Với mọi $1\le i<n$, $1\le j<m$, bốn ô
  $(i,j)$, $(i,j+1)$, $(i+1,j)$ và $(i+1,j+1)$ không thể có cùng màu.
\end{enumerate}

Một số ô đã được tô màu; bạn cần tô các ô còn lại.

Hãy in số phương án tô màu tốt modulo $998244353$.

\subsection*{Input}

Dòng đầu tiên của dữ liệu vào chứa một số nguyên $T$ ($T\le 10$), là số bộ dữ
liệu.

Mỗi bộ dữ liệu chứa $n+1$ dòng.

Dòng đầu tiên chứa hai số nguyên $n,m$
($3\le n\le 100$, $3\le m\le 100$, $nm\le 100$), là kích thước của lưới.

$n$ dòng tiếp theo mô tả các ô đã tô màu, mỗi dòng chứa $m$ số nguyên.

Số thứ $j$ trong hàng thứ $i$ mô tả ô $(i,j)$; số đó là $0$, $1$ hoặc $-1$.
Giá trị $0$ nghĩa là ô được tô trắng, $1$ nghĩa là ô được tô đen, và $-1$
nghĩa là ô chưa được tô.

\subsection*{Output}

In số phương án tô màu tốt modulo $998244353$.

\subsection*{Example}

\begin{verbatim}
standard input
3
3 3
1 0 0
1 1 0
-1 -1 0
3 4
1 -1 -1 0
-1 -1 -1 -1
0 -1 -1 1
10 10
-1 -1 -1 -1 -1 -1 -1 -1 -1 -1
-1 -1 -1 -1 -1 -1 -1 -1 -1 -1
-1 -1 -1 -1 -1 -1 -1 -1 -1 -1
-1 -1 -1 -1 -1 -1 -1 -1 -1 -1
-1 -1 -1 -1 -1 -1 -1 -1 -1 -1
-1 -1 -1 -1 -1 -1 -1 -1 -1 -1
-1 -1 -1 -1 -1 -1 -1 -1 -1 -1
-1 -1 -1 -1 -1 -1 -1 -1 -1 -1
-1 -1 -1 -1 -1 -1 -1 -1 -1 -1
-1 -1 -1 -1 -1 -1 -1 -1 -1 -1

standard output
2
0
139719073
\end{verbatim}

\end{document}
